\documentclass[12pt, a4paper, oneside]{book}

% --- Paquetes Requeridos ---
\usepackage[utf8]{inputenc}
\usepackage[spanish]{babel}
\usepackage{xcolor}
\usepackage{geometry}
\usepackage{tcolorbox}
\usepackage{titlesec}
\usepackage{fancyhdr}
\usepackage{graphicx}
\usepackage{hyperref}
\usepackage{enumitem}
\usepackage{listings}
\usepackage{amsmath}

% --- Configuración de Página ---
\geometry{top=2.5cm, bottom=2.5cm, left=3cm, right=2.5cm}

% --- Colores Ganesha (Premium Dark Theme) ---
\definecolor{brandorange}{HTML}{EA580C}
\definecolor{brandblack}{HTML}{09090B}
\definecolor{brandgray}{HTML}{27272A}
\definecolor{brandwhite}{HTML}{FAFAFA}

% --- Estilos de Títulos ---
\titleformat{\chapter}[display]{\normalfont\huge\bfseries\color{brandorange}}{\chaptertitlename\ \thechapter}{20pt}{\Huge}
\titleformat{\section}{\color{brandwhite}\normalfont\Large\bfseries}{\thesection}{1em}{}

% --- Configuración de Código ---
\lstset{
    backgroundcolor=\color{brandgray},
    basicstyle=\ttfamily\small\color{brandwhite},
    keywordstyle=\color{brandorange},
    stringstyle=\color{green},
    breaklines=true,
    frame=single,
    rulecolor=\color{brandorange}
}

\begin{document}

\pagecolor{brandblack}
\color{brandwhite}

% ==========================================
% PORTADA
% ==========================================
\begin{titlepage}
    \centering
    \vspace*{2cm}
    {\fontsize{50}{60}\selectfont\bfseries\color{brandorange} GANESHA}\\[0.5cm]
    {\Huge\bfseries ENTERPRISE ERP}\\[2cm]
    
    \begin{tcolorbox}[colback=brandgray, colframe=brandorange, arc=10pt]
        \centering
        \Large\textbf{Manual Maestro de Operaciones y Arquitectura Técnica}
    \end{tcolorbox}
    
    {\Large \textbf{Arquitecto: Wendell Flashey}}\\[0.5cm]
    {\Large \textbf{Versión Final de Auditoría (v4.0 Maestro Edition)}}\\[0.5cm]
    {\large \today}
\end{titlepage}

\tableofcontents
\newpage

% ==========================================
\chapter{Introducción y Visión}
% ==========================================
Ganesha Elite es una plataforma de gestión de recursos empresariales diseñada bajo la visión de Wendell Flashey para la era de la Inteligencia Artificial Soberana. A diferencia de los sistemas tradicionales, Ganesha integra un motor de auditoría forense en tiempo real que previene el fraude y garantiza la integridad contable absoluta, permitiendo que la tecnología hable el idioma del empresario.

\section{Filosofía de Diseño}
El sistema se rige por tres pilares:
\begin{itemize}
    \item \textbf{Inmutabilidad}: Una vez cerrado un período, los datos son inalterables.
    \item \textbf{Trazabilidad}: Cada clic es auditado con el actor y el cambio realizado.
    \item \textbf{Inteligencia}: Análisis predictivo de flujo de caja y tendencias de gasto.
\end{itemize}

% ==========================================
\chapter{Arquitectura de Seguridad}
% ==========================================
En esta versión 2.5, se ha implementado un esquema de seguridad de grado bancario para proteger la información multi-empresa.

\section{Firma de Tokens HMAC SHA256}
Toda comunicación entre el cliente y el servidor está firmada. Esto garantiza que el token de sesión no pueda ser falsificado ni reutilizado en otros dominios.

\section{Protección de Períodos Contables}
El sistema implementa el helper \texttt{ensurePeriodOpen}. Este mecanismo bloquea físicamente las operaciones de \texttt{POST, PUT y DELETE} si la fecha de la transacción pertenece a un mes fiscal cerrado.

% ==========================================
\chapter{Módulos Centrales}
% ==========================================

\section{Contabilidad y Partida Doble}
El motor contable valida matemáticamente que:
\begin{equation}
    \sum \text{Débitos} = \sum \text{Créditos}
\end{equation}
Antes de permitir que un asiento pase del estado \texttt{DRAFT} a \texttt{POSTED}.

\section{Gestión de Activos y Depreciación}
Ganesha automatiza el ciclo de vida del patrimonio. Al registrar un activo fijo (computadoras, vehículos), el sistema calcula la depreciación mensual lineal y genera el asiento contable de gasto de forma autónoma cada cierre de mes.

\section{Bancos y Tesorería}
El módulo de conciliación permite cruzar los movimientos bancarios reales con las pólizas de pago generadas desde las facturas, permitiendo una visión de liquidez en tiempo real.

% ==========================================
\chapter{Ganesha AI: Auditoría Forense}
% ==========================================
La Inteligencia Artificial integrada (Llama 3.2) no es solo un chatbot; es una analista financiera.
\begin{itemize}
    \item \textbf{Detección de Fugas}: Puede identificar gastos inusualmente altos en categorías específicas.
    \item \textbf{Análisis de Riesgo}: Evalúa la cartera vencida y predice problemas de liquidez a 30 días.
\end{itemize}

% ==========================================
\chapter{Guía del Desarrollador}
% ==========================================
El backend utiliza Next.js con Prisma ORM y SQLite como motor de persistencia optimizado.

\section{Esquema de Base de Datos}
El esquema cuenta con 22 modelos relacionales que garantizan la multi-tenencia mediante el campo \texttt{companyId} en todas las entidades críticas.

\begin{lstlisting}[language=SQL]
-- Ejemplo de integridad en Auditoria
CREATE TABLE AuditLog (
  id String PRIMARY KEY,
  userId String,
  action String, -- CREATE, UPDATE, DELETE
  oldValues JSON,
  newValues JSON
);
\end{lstlisting}

\vfill
\centering
\textbf{Fin del Manual Maestro}
\end{document}
